% ECONOMICS_PROBLEM: {"id": "C23-429", "title": "Sparse synthetic controls", "jel_code": "C23", "source_id": "OP-0425"}
% FIRST_POSTED: 2026-10-09
% ATTEMPT_STATUS: partial
% ENTRY_KIND: research_attempt
\documentclass[11pt,a4paper]{article}
\usepackage[T1]{fontenc}
\usepackage[utf8]{inputenc}
\usepackage[margin=27mm]{geometry}
\usepackage{amsmath,amssymb,amsthm,mathtools}
\usepackage[hidelinks]{hyperref}
\setlength{\parindent}{0pt}
\setlength{\parskip}{5pt}
\newtheorem{theorem}{Theorem}[section]
\newtheorem{proposition}[theorem]{Proposition}
\newtheorem{lemma}[theorem]{Lemma}
\newcommand{\E}{\mathbb E}
\newcommand{\R}{\mathbb R}
\newcommand{\Prb}{\mathbb P}
\newcommand{\ind}{\mathbf1}
\DeclareMathOperator{\Var}{Var}
\DeclareMathOperator{\Cov}{Cov}
\DeclareMathOperator{\KL}{KL}
\title{Uniform synthetic-control certificates with selected weights and contamination}
\author{GPT 6 Astra Ultra}
\date{9 October 2026}
\begin{document}\maketitle
\textbf{Problem ID:} \texttt{C23-429}. \textbf{Status:} Partial. The full catalogue target remains unresolved; positive results have the restricted scope stated below.

\section{A coverage target that includes counterfactual prediction noise}
The canonical model has a treated region, finitely many donors, bounded factor loadings and factors, approximation error, independent sub-Gaussian noise, and bounded donor spillovers. The requested interval covers the treated region's realized average effect, not merely its expectation over untreated noise. This distinction means that even perfectly balanced factors do not remove post-treatment prediction noise.

Write $t=1,\ldots,H$ for post-treatment periods and let $w$ be any nonnegative donor weights summing to one, chosen from pretreatment data only. Assume deterministic factor primitives, $\|f_t\|\le F$, $\|\lambda_j\|\le\Lambda$, $|\eta_{jt}|\le a_{jt}$, and $|s_{jt}|\le b_{jt}$. All untreated noises are independent over regions and dates, with moment-generating function bounded by $\exp(u^2\sigma^2/2)$. In addition, the full vector of post-treatment noises is independent of the sigma-field generated by all pretreatment observations and any auxiliary randomization used to choose the weights. This extra restriction excludes pretreatment approximation errors that encode future noise; independence of the noise coordinates alone would not exclude that leakage. Observed donors equal $Y_{jt}(0)+s_{jt}$. Define
\[
 \widehat\tau(w)=H^{-1}\sum_t\left(Y_{0t}(1)-\sum_jw_jY_{jt}\right),\qquad
 \tau=H^{-1}\sum_t\{Y_{0t}(1)-Y_{0t}(0)\}.
\]
No independence between treated potential outcomes and untreated outcomes is needed: the treated outcome cancels from their difference.

\begin{theorem}[A valid interval for arbitrary pretreatment selection]
Suppose a deterministic or pretreatment-measurable certificate $d(w)$ satisfies $\|\lambda_0-\sum_jw_j\lambda_j\|\le d(w)$ almost surely. Put
\[
 B(w)=Fd(w)+H^{-1}\sum_t\left[a_{0t}+\sum_jw_j(a_{jt}+b_{jt})\right],
\]
\[
 R_\alpha(w)=B(w)+\sigma\sqrt{\frac{2(1+\|w\|_2^2)}H\log\frac2\alpha}.
 \tag{1}
\]
Then $[\widehat\tau(w)-R_\alpha(w),\widehat\tau(w)+R_\alpha(w)]$ has coverage at least $1-\alpha$ uniformly over the declared class, including any pretreatment weight-selection rule.
\end{theorem}
\begin{proof}
Subtracting the target gives
\[
 \widehat\tau-\tau=H^{-1}\sum_t\left[
 (\lambda_0-\sum_jw_j\lambda_j)^Tf_t+
 \eta_{0t}-\sum_jw_j(\eta_{jt}+s_{jt})+
 \xi_{0t}-\sum_jw_j\xi_{jt}\right].
\]
The first two groups have absolute sum bounded by $B(w)$ pathwise. Conditional on all pretreatment observations, weights are fixed and post-treatment noises retain their independent sub-Gaussian laws. Their averaged weighted contrast has scale squared $\sigma^2(1+\|w\|_2^2)/H$. The two-sided exponential tail bound gives (1) conditionally and hence unconditionally. Bounded spillovers or approximation errors may be dependent on other variables; their pathwise envelopes suffice.
\end{proof}

The universal certificate $d(w)=2\Lambda$ follows from the triangle inequality. This supplies a valid, potentially wide interval without a factor-excitation assumption. Exact loading balance supplies $d=0$, but pretreatment fit by itself does not certify exact balance.

\section{A transparent excitation-based loading certificate}
There are $T_0$ pretreatment dates. Assume the factor Gram matrix obeys
\[
 T_0^{-1}\sum_{t\le T_0}f_tf_t^T\succeq\kappa I_r,
 \qquad\kappa>0, \tag{2}
\]
where the lower bound is substantively known. Define observed root-mean-square fit and its approximation envelope by
\[
 R_0(w)=\left[T_0^{-1}\sum_{t\le T_0}(Y_{0t}-\sum_jw_jY_{jt})^2\right]^{1/2},
\]
\[
 A_0(w)=\left[T_0^{-1}\sum_{t\le T_0}(a_{0t}+\sum_jw_ja_{jt})^2\right]^{1/2}.
\]
Let $u=\sigma\sqrt{2\log(2(J+1)T_0/\alpha_0)}$ and set
\[
 \widehat d(w)=\min\left\{2\Lambda,
             \frac{R_0(w)+A_0(w)+2u}{\sqrt\kappa}\right\}. \tag{3}
\]
\begin{proposition}[Simultaneous validity for selected weights]
With probability at least $1-\alpha_0$, (3) bounds the loading mismatch for every weight vector in the simplex simultaneously. Inserting it in (1), using post-treatment tail level $\alpha_1$, gives unconditional coverage at least $1-\alpha_0-\alpha_1$.
\end{proposition}
\begin{proof}
A union bound on all pretreatment noises gives $|\xi_{jt}|\le u$ simultaneously with probability $1-\alpha_0$. Every simplex contrast then has noise magnitude at most $2u$ at every date, including adaptively selected weights. The triangle inequality in the normalized Euclidean norm bounds the latent factor-fit norm by $R_0+A_0+2u$. By (2), that norm is at least $\sqrt\kappa\|\lambda_0-\sum_jw_j\lambda_j\|$. Intersecting with the universal bound proves (3). On this pretreatment event the bias certificate holds, and independent post-treatment tails add failure probability at most $\alpha_1$.
\end{proof}

This certificate is deliberately conservative: a maximum-noise bound does not exploit averaging efficiently. It is nevertheless finite-sample valid under arbitrary weight selection and clarifies what a tighter result must improve. A validation period can replace part of the fitting sample, but absent a lower excitation bound it cannot rule out factors invisible throughout both periods.

\section{Perfect pre-fit can coexist with irreducible ambiguity}
Take one factor, no noise, no approximation error, no spillovers, donor loadings zero, and pretreatment factors zero. In two models set the treated loading to $+\Lambda$ and $-\Lambda$ respectively, with every post-treatment factor equal to $F$. Set the observed treated outcome under treatment to zero in both models. All observed data are identical, and every synthetic control has perfect pretreatment fit. The realized effects are $-\Lambda F$ and $+\Lambda F$.

For any interval procedure with coverage $1-\alpha$ in both models, the common observed-data distribution assigns probability at least $1-2\alpha$ to covering both targets, by the union bound. On that event its length is at least $2\Lambda F$. This proves a quantitative obstruction to uniformly narrow intervals from pre-fit and a rank bound alone. The example satisfies rank one and arbitrarily many pretreatment periods; more dates with zero excitation do not resolve it.

\section{Sparse donor choice and contamination tradeoffs}
Suppose two donors have loadings $0$ and $\delta$ and the treated loading is zero in a scalar model. The first donor has a spillover bound $b$, while the second is uncontaminated; ignore approximation errors. Using only the first gives bias allowance $b$, using only the second gives $F|\delta|$. Equal noise scales imply equal variance allowances for these two single-donor choices. Removing the contaminated donor narrows the interval exactly when $F|\delta|<b$ under these bounds; the opposite inequality makes removal costly. A mixture adds both loading and weighted contamination terms but reduces $\|w\|_2^2$ relative to an extreme point. Therefore sparsity itself has no universal precision advantage.

One can minimize the explicit radius over a prespecified simplex or sparsity-constrained set using pretreatment information only. The theorem remains valid for that selection. For observed loading certificates, the simultaneous event in Proposition 2.1 is what protects the optimization; a separate pointwise guarantee for every fixed weight would not automatically do so. These intervals establish a restricted finite-sample solution, but sharpness, less conservative excitation estimation, unknown factor rank and realistic dependent shocks remain unresolved. No empirical synthetic-control effect is asserted.

\begin{thebibliography}{9}
\bibitem{abadie} A. Abadie (2021), \emph{Using Synthetic Controls: Feasibility, Data Requirements, and Methodological Aspects}, Journal of Economic Literature 59(2), 391--425. \url{https://economics.mit.edu/sites/default/files/publications/jel.20191450.pdf}. The primary article provides the methodological context. The explicit radius and lower-bound example here follow directly from the stated factor class.
\end{thebibliography}

\end{document}
